Настоящая работа посвящена исследованию и практической реализации фундаментальных алгоритмов на графах. Граф~--- это абстрактный математический объект, представляющий собой множество вершин и множество соединяющих их рёбер. В~современных информационных системах ориентированные графы используются для моделирования широкого спектра прикладных задач, включая транспортную логистику, социальные сети и сетевую маршрутизацию. Актуальность исследования обусловлена необходимостью эффективного поиска оптимальных путей в ориентированных структурах данных, что требует детального теоретического и практического анализа существующих подходов.

Поиск кратчайших путей и обход структуры графа~--- базовые операции для большинства сетевых алгоритмов. В~зависимости от специфики задачи, например, наличия рёбер с~отрицательными весами или необходимости работы в~режиме реального времени, выбираются различные алгоритмические решения. Ошибка при выборе структуры данных или алгоритма может привести к~критическому падению производительности программного комплекса, поэтому разработчику важно понимать границы применимости каждого метода.

Целью данной работы является систематизация теоретических сведений об основных алгоритмах на графах, сравнительный анализ их вычислительной сложности, а~также разработка и тестирование практической реализации алгоритмов на~языке программирования C++.
